% Source: https://artofproblemsolving.com/community/c3h1295259_its_the_final_problem
% Stable identity: aops-topic-1295259-q01
Put the rectangle at $(0,0),(8,0),(8,7),(0,7)$. The lower-left triangle has legs from $(0,0)$ to $(4,0)$ and to $(0,3)$, so its hypotenuse lies on the line $3x+4y=12$. The upper-left triangle has legs from $(0,7)$ to $(3,7)$ and down to $(0,3)$, so its hypotenuse joins $(0,3)$ and $(3,7)$ and lies on the line $4x-3y+9=0$. The figure is symmetric about the line $x=4$, and so are the two remaining hypotenuses, so we look for the circle with centre $O=(4,k)$ on this line.

\begin{asy}
usepackage("amsmath");
size(7cm);
defaultpen(fontsize(10pt));
pair A0=(0,0),B0=(8,0),C0=(8,7),D0=(0,7);
draw(A0--B0--C0--D0--cycle,linewidth(0.8));
pair P1=(4,0),P2=(0,3),P3=(0,3),P4=(3,7),P5=(8,3),P6=(5,7);
fill(A0--(4,0)--(0,3)--cycle,rgb(0.85,0.92,1)); fill(B0--(4,0)--(8,3)--cycle,rgb(0.85,0.92,1));
fill(D0--(3,7)--(0,3)--cycle,rgb(1,0.9,0.8)); fill(C0--(5,7)--(8,3)--cycle,rgb(1,0.9,0.8));
draw((4,0)--(0,3)); draw((4,0)--(8,3)); draw((3,7)--(0,3)); draw((5,7)--(8,3));
draw(A0--B0--C0--D0--cycle,linewidth(0.8));
pair O=(4,25/7); real r=20/7;
draw(circle(O,r),blue+linewidth(0.9)); dot(O,blue);
pair T1=(16/7,9/7), T2=(12/7,37/7);
draw(O--T1,dashed+blue); draw(O--T2,dashed+blue); dot(T1); dot(T2);
label("$O$",O,E,blue);
label("$4$",(2,0),S); label("$4$",(6,0),S); label("$3$",(0,1.5),W); label("$4$",(0,5),W);
label("$3$",(1.5,7),N); label("$3$",(6.5,7),N); label("$3$",(8,1.5),E); label("$4$",(8,5),E);
label("$\dfrac{20}{7}$",(O+T1)/2,NW,blue+fontsize(8pt)); label("$\dfrac{20}{7}$",(O+T2)/2,SW,blue+fontsize(8pt));
\end{asy}

The distances from $O$ to the two lines are
\[\dfrac{|3\cdot4+4k-12|}{5}=\dfrac{4k}{5}\qquad\text{and}\qquad\dfrac{|16-3k+9|}{5}=\dfrac{|25-3k|}{5}.\]
Tangency to all four hypotenuses (by symmetry, the right-hand ones give the same equations) requires these to be equal to the radius. With $O$ inside the rectangle between the lines, $4k=25-3k$, so $k=\dfrac{25}{7}$ and the radius is $\dfrac{4k}{5}=\dfrac{20}{7}$.

The points of tangency are the feet of the perpendiculars from $O$, namely $\left(\dfrac{16}{7},\dfrac{9}{7}\right)$ on the lower-left hypotenuse and $\left(\dfrac{12}{7},\dfrac{37}{7}\right)$ on the upper-left one; both lie on the hypotenuse segments (their $x$-coordinates are between $0$ and $4$, and between $0$ and $3$), so the circle really touches the four hypotenuses.

The diameter is $\dfrac{40}{7}$, so $m+n=40+7=\boxed{47}$.
